Recent work has explored optimizing string prompts, 
including gradient-guided search \citep{shin-etal-2020-autoprompt, wen2023hard}, reranking brute force search \citep{gao-etal-2021-making}, evolutionary algorithms \citep{fernando2023promptbreeder}, prompting other LMs \citep{yang2023large, zhou2023large, pryzant2023automatic}, and reinforcement learning (RL) \citep{deng-etal-2022-rlprompt, zhang2022tempera, hao2022optimizing}.  Prior work on RL for prompts focuses on word level edits of only a few words \citep{deng-etal-2022-rlprompt}, phrase level edits \citep{zhang2022tempera}, or text-to-image generation \citep{hao2022optimizing}. 
\citet{khattab2024dspy}  present DSPy, a programming model for expressing LM programs and optimizing their prompts 
and weights. Unlike our work, the authors only explore optimizers based on bootstrapping strong demonstrations. \citet{sordoni2023joint} explore joint prompt optimization for stacked LLM calls, modeling this as variational inference and exploring this for two simple layers. Their approach inherently relies on having access to log probabilities for explicitly passed tokens to LMs, an increasingly restrictive assumption in practice (e.g.\ with commodified LM APIs that allow only text-in-text-out processing). In contrast, our optimizers work on an arbitrary number of modules for any LM program out-of-the-box.
